% Options for packages loaded elsewhere
\PassOptionsToPackage{unicode}{hyperref}
\PassOptionsToPackage{hyphens}{url}
\documentclass[
]{article}
\usepackage{xcolor}
\usepackage[margin=1in]{geometry}
\usepackage{amsmath,amssymb}
\setcounter{secnumdepth}{-\maxdimen} % remove section numbering
\usepackage{iftex}
\ifPDFTeX
  \usepackage[T1]{fontenc}
  \usepackage[utf8]{inputenc}
  \usepackage{textcomp} % provide euro and other symbols
\else % if luatex or xetex
  \usepackage{unicode-math} % this also loads fontspec
  \defaultfontfeatures{Scale=MatchLowercase}
  \defaultfontfeatures[\rmfamily]{Ligatures=TeX,Scale=1}
\fi
\usepackage{lmodern}
\ifPDFTeX\else
  % xetex/luatex font selection
\fi
% Use upquote if available, for straight quotes in verbatim environments
\IfFileExists{upquote.sty}{\usepackage{upquote}}{}
\IfFileExists{microtype.sty}{% use microtype if available
  \usepackage[]{microtype}
  \UseMicrotypeSet[protrusion]{basicmath} % disable protrusion for tt fonts
}{}
\makeatletter
\@ifundefined{KOMAClassName}{% if non-KOMA class
  \IfFileExists{parskip.sty}{%
    \usepackage{parskip}
  }{% else
    \setlength{\parindent}{0pt}
    \setlength{\parskip}{6pt plus 2pt minus 1pt}}
}{% if KOMA class
  \KOMAoptions{parskip=half}}
\makeatother
\usepackage{color}
\usepackage{fancyvrb}
\newcommand{\VerbBar}{|}
\newcommand{\VERB}{\Verb[commandchars=\\\{\}]}
\DefineVerbatimEnvironment{Highlighting}{Verbatim}{commandchars=\\\{\}}
% Add ',fontsize=\small' for more characters per line
\usepackage{framed}
\definecolor{shadecolor}{RGB}{248,248,248}
\newenvironment{Shaded}{\begin{snugshade}}{\end{snugshade}}
\newcommand{\AlertTok}[1]{\textcolor[rgb]{0.94,0.16,0.16}{#1}}
\newcommand{\AnnotationTok}[1]{\textcolor[rgb]{0.56,0.35,0.01}{\textbf{\textit{#1}}}}
\newcommand{\AttributeTok}[1]{\textcolor[rgb]{0.13,0.29,0.53}{#1}}
\newcommand{\BaseNTok}[1]{\textcolor[rgb]{0.00,0.00,0.81}{#1}}
\newcommand{\BuiltInTok}[1]{#1}
\newcommand{\CharTok}[1]{\textcolor[rgb]{0.31,0.60,0.02}{#1}}
\newcommand{\CommentTok}[1]{\textcolor[rgb]{0.56,0.35,0.01}{\textit{#1}}}
\newcommand{\CommentVarTok}[1]{\textcolor[rgb]{0.56,0.35,0.01}{\textbf{\textit{#1}}}}
\newcommand{\ConstantTok}[1]{\textcolor[rgb]{0.56,0.35,0.01}{#1}}
\newcommand{\ControlFlowTok}[1]{\textcolor[rgb]{0.13,0.29,0.53}{\textbf{#1}}}
\newcommand{\DataTypeTok}[1]{\textcolor[rgb]{0.13,0.29,0.53}{#1}}
\newcommand{\DecValTok}[1]{\textcolor[rgb]{0.00,0.00,0.81}{#1}}
\newcommand{\DocumentationTok}[1]{\textcolor[rgb]{0.56,0.35,0.01}{\textbf{\textit{#1}}}}
\newcommand{\ErrorTok}[1]{\textcolor[rgb]{0.64,0.00,0.00}{\textbf{#1}}}
\newcommand{\ExtensionTok}[1]{#1}
\newcommand{\FloatTok}[1]{\textcolor[rgb]{0.00,0.00,0.81}{#1}}
\newcommand{\FunctionTok}[1]{\textcolor[rgb]{0.13,0.29,0.53}{\textbf{#1}}}
\newcommand{\ImportTok}[1]{#1}
\newcommand{\InformationTok}[1]{\textcolor[rgb]{0.56,0.35,0.01}{\textbf{\textit{#1}}}}
\newcommand{\KeywordTok}[1]{\textcolor[rgb]{0.13,0.29,0.53}{\textbf{#1}}}
\newcommand{\NormalTok}[1]{#1}
\newcommand{\OperatorTok}[1]{\textcolor[rgb]{0.81,0.36,0.00}{\textbf{#1}}}
\newcommand{\OtherTok}[1]{\textcolor[rgb]{0.56,0.35,0.01}{#1}}
\newcommand{\PreprocessorTok}[1]{\textcolor[rgb]{0.56,0.35,0.01}{\textit{#1}}}
\newcommand{\RegionMarkerTok}[1]{#1}
\newcommand{\SpecialCharTok}[1]{\textcolor[rgb]{0.81,0.36,0.00}{\textbf{#1}}}
\newcommand{\SpecialStringTok}[1]{\textcolor[rgb]{0.31,0.60,0.02}{#1}}
\newcommand{\StringTok}[1]{\textcolor[rgb]{0.31,0.60,0.02}{#1}}
\newcommand{\VariableTok}[1]{\textcolor[rgb]{0.00,0.00,0.00}{#1}}
\newcommand{\VerbatimStringTok}[1]{\textcolor[rgb]{0.31,0.60,0.02}{#1}}
\newcommand{\WarningTok}[1]{\textcolor[rgb]{0.56,0.35,0.01}{\textbf{\textit{#1}}}}
\usepackage{graphicx}
\makeatletter
\newsavebox\pandoc@box
\newcommand*\pandocbounded[1]{% scales image to fit in text height/width
  \sbox\pandoc@box{#1}%
  \Gscale@div\@tempa{\textheight}{\dimexpr\ht\pandoc@box+\dp\pandoc@box\relax}%
  \Gscale@div\@tempb{\linewidth}{\wd\pandoc@box}%
  \ifdim\@tempb\p@<\@tempa\p@\let\@tempa\@tempb\fi% select the smaller of both
  \ifdim\@tempa\p@<\p@\scalebox{\@tempa}{\usebox\pandoc@box}%
  \else\usebox{\pandoc@box}%
  \fi%
}
% Set default figure placement to htbp
\def\fps@figure{htbp}
\makeatother
\setlength{\emergencystretch}{3em} % prevent overfull lines
\providecommand{\tightlist}{%
  \setlength{\itemsep}{0pt}\setlength{\parskip}{0pt}}
\usepackage{bookmark}
\IfFileExists{xurl.sty}{\usepackage{xurl}}{} % add URL line breaks if available
\urlstyle{same}
\hypersetup{
  pdftitle={Highest Posterior Density Interval},
  hidelinks,
  pdfcreator={LaTeX via pandoc}}

\title{Highest Posterior Density Interval}
\author{}
\date{\vspace{-2.5em}2026-04-17}

\begin{document}
\maketitle

\subsection{Introduction}\label{introduction}

The highest posterior density (HPD) interval is the narrowest range that
contains a chosen probability mass --- typically 95 \% --- under a
posterior distribution. Two intervals can both cover 95 \% of the
posterior, yet the HPD will always be at least as short as the
equal-tailed interval (ETI) and will include only those parameter values
whose density is above some threshold. For unimodal posteriors the two
intervals are similar; for skewed or bimodal posteriors they can
disagree dramatically, and the choice between them shapes how readers
perceive the result.

\subsection{Prerequisites}\label{prerequisites}

A working understanding of posterior distributions and credible
intervals, and access to MCMC draws (via \texttt{brms},
\texttt{rstanarm}, raw Stan, or any other sampler).

\subsection{Theory}\label{theory}

For a posterior \(p(\theta \mid y)\), the \((1-\alpha)\) HPD region is

\[\{\theta : p(\theta \mid y) \ge h\}, \quad \text{with } h \text{ such that } P(\theta \in \text{region}) = 1 - \alpha.\]

For a unimodal posterior this region is a single interval centred near
the mode; for a multimodal posterior it can be a union of disjoint
intervals, one per mode that exceeds the threshold. The ETI, by
contrast, places the same tail probability \(\alpha/2\) in each tail, so
for a skewed posterior it shifts toward the heavier tail and is wider.
HPD therefore favours the shape of the posterior; ETI favours
interpretability and invariance to monotone transformations.

\subsection{Assumptions}\label{assumptions}

Adequate MCMC draws (after warm-up, with \(\hat R\) near 1.00) or a
closed-form posterior density. The HPD is sensitive to noisy density
estimates in the tails, so very short chains can yield unstable interval
endpoints.

\subsection{R Implementation}\label{r-implementation}

\begin{Shaded}
\begin{Highlighting}[]
\FunctionTok{library}\NormalTok{(HDInterval)}

\FunctionTok{set.seed}\NormalTok{(}\DecValTok{2026}\NormalTok{)}
\NormalTok{theta }\OtherTok{\textless{}{-}} \FunctionTok{c}\NormalTok{(}\FunctionTok{rbeta}\NormalTok{(}\DecValTok{5000}\NormalTok{, }\DecValTok{5}\NormalTok{, }\DecValTok{15}\NormalTok{), }\FunctionTok{rbeta}\NormalTok{(}\DecValTok{5000}\NormalTok{, }\DecValTok{15}\NormalTok{, }\DecValTok{5}\NormalTok{))  }\CommentTok{\# bimodal}

\CommentTok{\# Equal{-}tailed (may span both modes)}
\FunctionTok{quantile}\NormalTok{(theta, }\FunctionTok{c}\NormalTok{(}\FloatTok{0.025}\NormalTok{, }\FloatTok{0.975}\NormalTok{))}

\CommentTok{\# HPD (disjoint regions)}
\FunctionTok{hdi}\NormalTok{(theta, }\AttributeTok{credMass =} \FloatTok{0.95}\NormalTok{, }\AttributeTok{allowSplit =} \ConstantTok{TRUE}\NormalTok{)}
\end{Highlighting}
\end{Shaded}

\subsection{Output \& Results}\label{output-results}

\texttt{hdi()} returns a numeric vector of two endpoints when the HPD is
a single interval, or a matrix of begin / end pairs when
\texttt{allowSplit\ =\ TRUE} and the posterior is multimodal.
\texttt{quantile()} returns the ETI for direct comparison; the gap
between ETI and HPD widths is the most direct numerical signal of
posterior asymmetry.

\subsection{Interpretation}\label{interpretation}

A reporting sentence: ``The 95 \% HPD interval for \(\theta\) was
\((0.13, 0.38) \cup (0.60, 0.89)\), reflecting the bimodal posterior;
the equal-tailed interval \((0.13, 0.89)\) obscures this structure by
including the low-density valley between modes.'' When the posterior is
unimodal and roughly symmetric, the two intervals will agree and the
choice does not matter; when they disagree, prefer the one whose
definition matches the inferential question --- HPD for ``where does the
posterior concentrate?'' and ETI for ``what range of values is at most
2.5 \% above / below the truth?''

\subsection{Practical Tips}\label{practical-tips}

\begin{itemize}
\tightlist
\item
  Set \texttt{allowSplit\ =\ TRUE} whenever multimodality is plausible;
  a single interval reported across modes can be misleading.
\item
  HPD is \emph{not} invariant under non-linear reparameterisation: the
  HPD for \(\sigma^2\) is not in general the square of the HPD for
  \(\sigma\), while the ETI quantile-based interval is. Decide on the
  reporting scale before computing the interval.
\item
  For bounded parameters (probabilities, variances) the HPD often
  touches the boundary; report this explicitly so readers know the
  posterior is not symmetric there.
\item
  Use post-warm-up MCMC draws and check that the chains have mixed; HPD
  endpoints are noisier than means, so bigger effective sample sizes
  help.
\item
  \texttt{bayestestR::hdi()} provides a similar interface and integrates
  with model objects from \texttt{brms}, \texttt{rstanarm}, and
  \texttt{posterior}.
\end{itemize}

\subsection{R Packages Used}\label{r-packages-used}

\texttt{HDInterval} for the canonical \texttt{hdi()} function with
\texttt{allowSplit}, \texttt{bayestestR} for an integrated wrapper
around model objects, and \texttt{posterior} for handling MCMC draw
arrays consistently across backends.

\subsection{For Reviewers}\label{for-reviewers}

\textbf{What to look for in a paper using this method.}

\begin{itemize}
\tightlist
\item
  \textbf{Common misapplications.}
\item
  \textbf{Diagnostics that should be reported but often aren't.}
\item
  \textbf{Red flags in tables and figures.}
\item
  \textbf{What to verify.}
\item
  \textbf{What an adequate Methods paragraph must contain.}
\end{itemize}

\subsection{See also --- labs in this
chapter}\label{see-also-labs-in-this-chapter}

\begin{itemize}
\tightlist
\item
  \href{labs/lab-bayesian-thinking-control-vs-decision-error.Rmd}{Bayesian
  thinking; α control vs decision error}
\item
  \href{labs/lab-brms-stan-loo-hierarchical-models.Rmd}{brms / Stan,
  LOO, hierarchical models}
\end{itemize}

\end{document}
